\section{Represent black-hole QNMs by complex scaling}
\label{sec:bh-complex-scaling}
% BEGIN CURATED SUBJECT INDEX
\index{complex scaling}
\index{quasinormal mode}
% END CURATED SUBJECT INDEX

\subsection{Complex deformation of the tortoise coordinate}
\label{subsec:bh-csm-map}
% BEGIN CURATED SUBJECT INDEX
\index{quasinormal mode}
\index{tortoise coordinate}
% END CURATED SUBJECT INDEX

Set
\begin{equation}
  r_*=\rme^{\rmi\theta}x,
  \qquad x\in\real,
  \qquad \theta>0 .
  \label{eq:bh-complex-scaling}
\end{equation}
The complex-scaled master operator is
\begin{equation}
  H_{\theta}
  =-\rme^{-2\rmi\theta}
  \frac{\rmd^2}{\rmd x^2}
  +V_{sl}\!\left(r(\rme^{\rmi\theta}x)\right) .
  \label{eq:bh-H-theta}
\end{equation}
The QNM equation becomes
\begin{equation}
  H_\theta\psi_{sl}^{\theta}
  =\omega^2\psi_{sl}^{\theta} .
  \label{eq:bh-csm-eigenproblem}
\end{equation}
For a mode with $\omega=|\omega|\rme^{-\rmi\varphi}$, both asymptotic waves in
Eq.~\eqref{eq:qnm-boundary-flat} decay on the deformed line when
\begin{equation}
  \theta>\varphi .
  \label{eq:bh-frequency-exposure}
\end{equation}
Equivalently, $E=\omega^2$ lies above the continuum ray rotated by
$-2\theta$.

The formula is simple, but its analytic implementation is not automatic.
The inverse map $r(r_*)$ is multivalued because the tortoise coordinate
contains logarithms.  A deformation must specify branches and avoid metric
singularities.  One may work directly with an analytically continued
$V(r_*)$, use exterior scaling where the asymptotic form is controlled, or
formulate the contour in the complex $r$ plane.  A numerical result without a
documented branch choice is not reproducible.

\subsection{Why the spectral parameter is
\texorpdfstring{$\omega^2$}{omega squared}}
\label{subsec:omega-square}
% BEGIN CURATED SUBJECT INDEX
\index{spectrum!point}
\index{complex scaling}
\index{quasinormal mode}
\index{horizon boundary condition}
\index{Kerr black hole}
% END CURATED SUBJECT INDEX

Complex scaling rotates the continuum of the Schr\"odinger operator in the
$E$ plane.  Because $E=\omega^2$, the map to the frequency plane is
two-to-one.  A ray with $\arg E=-2\theta$ corresponds to frequency rays with
$\arg\omega=-\theta$ and $\arg\omega=\pi-\theta$.  The retarded QNM branch is
selected by the time convention and the horizon/infinity boundary conditions.
Plotting eigenvalues directly in the $\omega$ plane without tracking the
square-root branch can misidentify time-reversed or negative-frequency
partners.

If the potential depends on $\omega$, as happens after some transformations
of the Kerr or coupled-channel equations, the problem is a nonlinear
eigenvalue problem rather than Eq.~\eqref{eq:bh-csm-eigenproblem}.  Linearizing
it by enlarging the state space is possible only for particular rational or
polynomial dependence.  This is one reason the extension from static
spherical backgrounds to Kerr is nontrivial.

\subsection{Basis construction and parity}
\label{subsec:bh-basis}
% BEGIN CURATED SUBJECT INDEX
\index{spectrum!point}
\index{coupled-channel equation}
\index{quasinormal mode}
% END CURATED SUBJECT INDEX

The effective barrier is centered in a finite range of $r_*$.  One may use
even and odd Gaussian functions as in Eq.~\eqref{eq:real-gaussian-basis}, even
though the black-hole potential itself need not have exact reflection
symmetry.  Both parity families together span general functions.  Polynomial
times Gaussian bases,
\begin{equation}
  \phi_{np}(x)=x^p\rme^{-\alpha_nx^2},
  \qquad \text{$p=0$, $1$, \dots, $p_{\max}$} ,
  \label{eq:polynomial-gaussian}
\end{equation}
and complex-range Gaussians improve oscillatory resolution.  Nonorthogonal
blocks can grow or decay exponentially under scaling, so block balancing and
overlap truncation are essential for coupled channels
\cite{Ogawa:2026veu,Ogawa:2026dSCSM}.

The output should be analyzed in $E=\omega^2$ first.  A candidate QNM is an
isolated eigenvalue stable under $\theta$, $N$, basis widths, integration
range, and branch conventions.  Only then should a square root be taken to
quote $\omega$.  Low-lying, weakly damped modes are usually much more stable
than high overtones because they lie farther from the rotated continuum and
are represented by less oscillatory scaled wave functions.

\subsection{Schwarzschild and extremal RN benchmarks}
\label{subsec:bh-csm-benchmarks}
% BEGIN CURATED SUBJECT INDEX
\index{pseudostate}
\index{Regge--Wheeler equation}
\index{Schwarzschild black hole}
\index{Reissner--Nordstr\"om black hole}
% END CURATED SUBJECT INDEX

Leaver's continued-fraction method provides accurate benchmark frequencies
for Schwarzschild and Reissner--Nordstr\"om backgrounds
\cite{Leaver:1985ax,Leaver:1986gd,Nollert:1993zz}.  Complex-scaling
calculations reproduce low-lying Schwarzschild Regge--Wheeler modes and extend
to scalar, electromagnetic, and gravitational sectors of extremal
Reissner--Nordstr\"om
\cite{Ogawa:2026veu}.  The most robust agreement occurs for the fundamental
modes.  First overtones require broader basis scans, while still higher modes
can be difficult to distinguish from continuum pseudostates in an unoptimized
Gaussian basis.

This pattern has a spectral explanation.  For fixed $\theta$, the distance
between a pole and the rotated continuum controls isolation in the non-normal
eigenproblem.  Higher overtones have larger negative imaginary parts and lie
closer to, or below, the accessible ray.  Increasing $\theta$ exposes them but
also probes a wider complex-coordinate sector where the potential and basis
may be less well conditioned.  There is therefore no monotonic rule that a
larger scaling angle is numerically better.

\subsection{Spectral instability and pseudospectra}
\label{subsec:bh-pseudospectrum}
% BEGIN CURATED SUBJECT INDEX
\index{spectrum!point}
\index{resolvent}
\index{pseudospectrum}
\index{resonance!residue}
\index{complex scaling}
\index{quasinormal mode}
\index{continuum response}
% END CURATED SUBJECT INDEX

Black-hole QNM operators can be highly non-normal.  Define the
$\epsilon$-pseudospectrum of the matrix pencil $(H^\theta,N)$ by
\begin{equation}
  \pseudospectrum(H^\theta,N)
  =\left\{z\in\complex:
  s_{\min}(H^\theta-zN)<\epsilon\right\},
  \label{eq:generalized-pseudospectrum}
\end{equation}
where $s_{\min}$ is the smallest singular value after a documented basis
normalization.  Large pseudospectral contours indicate that a small potential
or discretization perturbation can move eigenvalues far from their unperturbed
positions
\cite{Jaramillo:2020tuu,Yang:2024vor}.

Spectral instability does not imply identical instability of all time-domain
observables.  The resolvent norm, residues, source projection, and duration of
the observation determine whether a shifted pole has appreciable effect.
Greybody factors and suitably integrated response functions can be more
stable than individual high overtones
\cite{Rosato:2024arw}.  Complex scaling is useful here because it gives access
to both poles and the deformed continuum resolvent.
